% 方法章节
% 说明：
% 1. 本文件设计为被主文档 \input{} 进来使用。
% 2. 内容来源于 04_paper.md 的 Method 部分。
% 3. 这里只做 Markdown 到 LaTeX 的语法迁移与公式环境适配。

\section{Method}

\subsection{系统概述}

本系统是一个基于外部位姿前端的增量式 2DGS 场景重建后端。系统接收外部 LiDAR-Inertial-Camera 前端（如 R3Live）实时输出的同步数据流，每帧包含 RGB 图像、稀疏 LiDAR 点云和 6-DOF 相机位姿，以 2DGS surfel 为基础地图表示。整体流程如图 X 所示。

当新关键帧到达时，系统首先将稀疏 LiDAR 点投影到图像平面形成稀疏深度图，经 SPNet 深度补全后在 LiDAR 盲区生成补充初始 3D 点；与此同时，DA3 单目深度估计模型预测稠密深度并与 LiDAR 对齐，由此为初始点提供法线先验。在初始点生成后，系统对当前高斯地图进行渲染，依据渲染透明度和颜色误差识别欠重建区域，并结合颜色梯度过滤与体素下采样，从中筛选出确需新增高斯的候选点子集。筛选后的候选点及其几何先验送入前馈回归网络，该网络利用共享 backbone 提取的当前帧与历史帧的原始与渲染图像特征，通过跨注意力机制融合多视角上下文，在单次前向传播中直接回归每个候选点的完整 surfel 属性。预测的新高斯插入地图后，系统在滑窗内的关键帧上进行多轮在线优化，以光度损失（L1 + SSIM）与 2DGS 几何约束损失（深度畸变 + 法线一致性）联合驱动高斯属性收敛。非关键帧仅接收位姿而不触发地图扩展。

前馈回归网络的训练通过自蒸馏闭环完成。系统首先以 \texttt{naive} 模式（不使用前馈网络，采用传统的基于规则初始化）在目标场景上运行一遍完整的增量式建图，收集经在线优化和离线精细化后的高斯属性作为伪真值，回溯构造训练样本，离线训练前馈网络后重新部署。

\subsection{2DGS 表示}

本系统采用 2D Gaussian Splatting（2DGS）作为地图基础表示。与标准 3DGS 的三维椭球原语不同，2DGS 将高斯压缩为嵌入三维空间的二维有向高斯盘（surfel），赋予每个原语明确的局部平面几何意义。每个 surfel 由中心位置 $\boldsymbol{\mu} \in \mathbb{R}^3$、两个正交切线方向 $\mathbf{t}_u$、$\mathbf{t}_v$ 及对应的二维尺度因子 $(s_u, s_v)$ 定义，法线方向由 $\mathbf{n} = \mathbf{t}_u \times \mathbf{t}_v$ 给出。方向信息以单位四元数 $\mathbf{q} \in \mathbb{S}^3$ 存储，其对应旋转矩阵 $\mathbf{R} = [\mathbf{t}_u \mid \mathbf{t}_v \mid \mathbf{n}]$。此外，每个 surfel 携带不透明度 $\alpha$ 和球谐颜色系数以表达视角依赖的外观。渲染阶段，像素颜色通过 alpha blending 由前到后逐层累积：

\begin{equation}
\mathbf{c}(\mathbf{x}) = \sum_{i} \mathbf{c}_i \, \alpha_i \, \hat{\mathcal{G}}_i(\mathbf{u}(\mathbf{x})) \prod_{j=1}^{i-1} \left(1 - \alpha_j \, \hat{\mathcal{G}}_j(\mathbf{u}(\mathbf{x}))\right)
\end{equation}

其中 $\hat{\mathcal{G}}_i$ 为 ray-splat intersection 计算的高斯权重。类似地，可对深度和法线进行相同的加权累积渲染几何信息。为充分利用 surfel 的表面化结构，系统在在线优化中引入深度畸变损失 $\mathcal{L}_{\text{dist}}$ 和法线一致性损失 $\mathcal{L}_{\text{normal}}$，分别促使沿射线的权重集中于单表面、并约束 surfel 法线与渲染深度推导的表面法线一致。

2DGS 的核心优势在于每个 surfel 天然具备明确的法线语义，这使得外部法线先验能够直接与 surfel 旋转参数耦合，为后续前馈旋转预测提供几何可解释的锚点。

\subsection{稠密几何先验}

为在增量式建图中生成足够密集且几何合理的初始高斯点，系统联合使用 SPNet 深度补全与 DA3 法线引导两个互补模块。

\subsubsection{SPNet 深度补全}

稀疏 LiDAR 点云在图像平面上覆盖有限，难以为 LiDAR 视场之外或遮挡区域提供几何观测。系统遵循 Gaussian-LIC2 的方案，采用 SPNet 轻量级深度补全网络将 RGB 图像与稀疏 LiDAR 深度融合为稠密深度图，并在无 LiDAR 覆盖的区域采样补充 3D 初始点。最终，SPNet 补充点与原始 LiDAR 投影点合并构成初始点集。

\subsubsection{DA3 法线引导}

系统使用 DA3（基于 Depth Anything v3 的单目深度估计模型）为每个初始点对应的 surfel 提供初始法线方向。DA3 对当前帧 RGB 图像推理输出相对深度图后，系统通过最小二乘仿射对齐将其与稀疏 LiDAR 深度配准，得到与真实尺度一致的稠密深度图，再以中心差分计算逐像素法线。按初始点的像素坐标在法线图上采样，得到每个初始点的相机空间法线 $\mathbf{n}_{\text{cam}} \in \mathbb{R}^3$。

需要指出的是，DA3 法线计算前置于前馈回归网络——法线先验首先由 DA3 独立计算完成，随后作为输入传入前馈网络，网络在此基础上进行微量修正和完整旋转构造（详见 III-D.4 节）。

\subsection{前馈高斯回归网络}

本节详细描述前馈网络的架构设计。该网络以候选点及其上游几何先验为输入，在单次前向传播中预测完整的 2DGS surfel 属性。

\subsubsection{逐点上下文构造}

前馈网络的输入准备核心在于逐点对齐的上下文构造（point-aligned context construction）。与 pixelSplat、MVSplat 等需要从图像中同时推断几何与高斯属性的离线纯视觉方法不同，本文已经通过 LiDAR 与深度先验获得候选点，因此无需再让网络隐式恢复位置；系统只需将候选点显式投影到共视窗口，并在逐点层面构造多视角上下文特征，使每个输入向量都与一个待插入的 surfel 一一对应。

特征提取模块沿用 pixelSplat 的 backbone 设计，从多尺度 ResNet 特征层经 $1 \times 1$ 卷积投影到统一维度，再经双线性上采样后求和，得到与输入图像同分辨率的逐像素特征图。为降低训练成本并提高稳定性，系统直接复用 pixelSplat 发布权重中的 backbone 参数，并在训练过程中保持冻结。该共享 backbone 对四类图像分别提取特征图：当前帧原始图像 $I_t \to F_t^{\text{img}}$，当前帧渲染图像 $\hat{I}_t \to F_t^{\text{ren}}$（用当前高斯地图从当前视角渲染），滑窗内历史帧原始图像 $\{I_{t-k}\} \to \{F_{t-k}^{\text{img}}\}$，以及历史帧渲染图像 $\{\hat{I}_{t-k}\} \to \{F_{t-k}^{\text{ren}}\}$（用当前高斯地图从历史视角渲染）。渲染图像的引入使网络能感知当前地图在该位置的重建状态，渲染特征与原始特征的差异隐式指示了该区域的重建不足程度。

在当前帧特征采样阶段，对每个候选点，以其像素坐标 $\mathbf{u}_i$ 在 $F_t^{\text{img}}$ 和 $F_t^{\text{ren}}$ 上进行双线性插值采样，得到逐点特征 $\mathbf{f}_i^{\text{curr}} \in \mathbb{R}^{D}$ 和 $\mathbf{f}_i^{\text{ren}} \in \mathbb{R}^{D}$，同时计算当前帧下的归一化观察方向 $\mathbf{d}_i^{\text{curr}} = \text{normalize}(R_{\text{cw}} \boldsymbol{\mu}_i + \mathbf{t}_{\text{cw}})$。

在历史帧投影与特征采样阶段，对共视窗口内 $W$ 帧历史关键帧，逐帧将候选点的世界坐标投影到该历史帧的图像平面，进行可见性判断（正向深度且落在图像范围内），在通过可见性的像素位置上分别从 $F_{t-k}^{\text{img}}$ 和 $F_{t-k}^{\text{ren}}$ 采样，得到逐点历史特征 $\mathbf{f}_i^{\text{hist}} \in \mathbb{R}^{W \times D}$、历史渲染特征 $\mathbf{f}_i^{\text{hist\_ren}} \in \mathbb{R}^{W \times D}$ 及可见性掩码 $\mathbf{m}_i^{\text{hist}} \in \{0,1\}^{W}$。同时计算每帧的观察方向（统一到当前帧坐标系表达）：$\mathbf{d}_{i,k}^{\text{hist}} = \text{normalize}(R_{\text{cw}}^t (\boldsymbol{\mu}_i - \mathbf{c}_{t-k}))$，其中 $\mathbf{c}_{t-k}$ 为历史帧相机中心。该方向编码随后与历史特征一同送入注意力模块。

\subsubsection{跨注意力时序上下文融合}

为将 $W$ 帧历史上下文聚合为固定维度的上下文向量，系统采用单头跨注意力模块。查询向量 $\mathbf{q}_i$ 由当前帧的观察方向、图像特征和渲染特征拼接后经线性投影得到；键和值向量 $\mathbf{k}_{i,k}$、$\mathbf{v}_{i,k}$ 则由对应历史帧的同类信息构造。上下文向量通过带可见性掩码的缩放点积注意力计算：

\begin{equation}
\mathbf{ctx}_i = \text{softmax}\!\left(\frac{\mathbf{q}_i \mathbf{K}_i^\top}{\sqrt{d_k}} + \mathbf{M}_i\right) \mathbf{V}_i
\end{equation}

其中 $d_k = 64$，$\mathbf{M}_i$ 将不可见帧的对应位置填充为 $-\infty$ 以完全排除其贡献。该设计的关键在于：观察方向显式参与查询和键的投影，使注意力权重能隐式编码视差关系，接近当前视角的历史帧自然获得更高权重。键和值共享相同的历史帧输入，使历史信息同时参与匹配和内容聚合。

\subsubsection{几何感知编码器}

跨注意力输出的上下文向量 $\mathbf{ctx}_i$ 与当前帧图像/渲染特征、逆深度 $\bar{z}_i^{-1}$、DA3 相机空间法线 $\mathbf{n}_i^{\text{cam}}$、有效性掩码 $m_i$ 和 KNN 像素间距 $d_i^{\text{nn}}$ 拼接后送入残差 MLP 编码器。其中逆深度、法线和像素间距均经 NeRF 风格的位置编码（频带数分别为 4、2、3），拼接后输入总维度为 544。编码器由一层 $544 \to 512$ 线性投影和 4 个 pre-activation 残差块组成，输出 512 维特征 $\mathbf{h}_i$。

\subsubsection{多头预测与几何约束解码}

编码器输出的 512 维特征向量分支到多个独立预测头。每个头具有 $\text{Linear} \to \text{ReLU} \to \text{Linear}$ 结构。关键设计在于：各头的输出不直接作为最终高斯参数，而是经过几何约束解码，将网络预测锚定在上游几何先验上。

\paragraph{(a) 法线引导的旋转构造}

系统不直接回归四元数（无约束空间中容易产生几何不一致的解），而是以 DA3 提供的相机空间法线 $\mathbf{n}_i^{\text{cam}}$ 为锚点，由网络预测修正量，通过正交化过程推导完整旋转。

网络首先预测法线修正向量 $\delta \mathbf{z}_i \in \mathbb{R}^3$、混合强度 $\lambda_i = \text{softplus}(\cdot) \in \mathbb{R}_{>0}$ 和辅助向量 $\mathbf{a}_i \in \mathbb{R}^3$，由此得到修正后的法线方向：

\begin{equation}
\mathbf{z}_i = \text{normalize}(\mathbf{n}_i^{\text{cam}} + \lambda_i \cdot \delta \mathbf{z}_i)
\end{equation}

随后通过 Gram-Schmidt 正交化过程，利用辅助向量与修正法线的叉积构造完整正交基：

\begin{equation}
\mathbf{y}_i = \text{normalize}(\mathbf{z}_i \times \mathbf{a}_i), \qquad \mathbf{x}_i = \mathbf{y}_i \times \mathbf{z}_i
\end{equation}

最终，相机系旋转矩阵 $[\mathbf{x}_i \mid \mathbf{y}_i \mid \mathbf{z}_i]$ 经当前帧位姿 $R_{\text{wc}}$ 转换到世界系并转为四元数存储。初始化设计确保网络起步时紧靠几何先验：$\lambda$ 头的 bias 初始化为 $-2.0$（使 $\text{softplus}(-2) \approx 0.13$，初始修正幅度很小），$\mathbf{a}$ 头的 bias 初始化为 $0.1$（避免退化叉积）。这保证了在训练初期，surfel 法线方向几乎完全由 DA3 先验决定，网络逐步学习必要的偏离。

\paragraph{(b) 密度自适应尺度预测}

网络预测 2D 尺度在像素平面上的归一化值，并通过局部点密度先验进行自适应缩放：

\begin{equation}
\mathbf{s}_i^{\text{pixel}} = \frac{d_i^{\text{nn}}}{2} \cdot \exp(\text{head}(\mathbf{h}_i)) \in \mathbb{R}^2_{>0}
\end{equation}

其中 $d_i^{\text{nn}}$ 为候选点的 KNN 像素间距。最终世界空间尺度通过 $\mathbf{s}_i^{\text{world}} = \mathbf{s}_i^{\text{pixel}} / (\bar{z}_i^{-1} \cdot f)$ 转换。这一参数化使网络无需学习深度依赖的绝对尺度，而只需学习与局部密度相关的相对尺度。

\paragraph{(c) 颜色组装}

DC 颜色头输出残差 $\Delta \mathbf{c}_i^{\text{dc}}$，加到由单帧 RGB 观测转换的基础 SH 先验上得到 $\mathbf{c}_i^{\text{dc}} = \text{RGB2SH}(\mathbf{I}_i) + \Delta \mathbf{c}_i^{\text{dc}}$；高阶 SH 系数由独立头直接预测。这种残差设计使网络从合理的颜色先验出发，仅需学习修正量。

\paragraph{(d) 不透明度}

不透明度头输出经 sigmoid 激活的真实不透明度 $\alpha_i = \sigma(\text{head}(\mathbf{h}_i))$，后续转为 logit 空间存储。

\subsection{自蒸馏训练闭环}

前馈网络的训练不依赖外部多视图标注数据集，而是通过以下自蒸馏管线完成。

\subsubsection{在线数据采集}

系统首先以 \texttt{naive} 模式（不使用前馈网络，采用传统的基于规则初始化）在目标场景上运行一次完整的增量式建图。运行过程中，每当某帧新增的高斯累积训练达到指定次数后，系统快照当前高斯属性作为中间伪真值；运行结束后再进行离线精细化（继续优化至收敛），得到最终高斯属性。随后，系统回溯每个关键帧，以因果正确的方式重新构造上下文特征：对每帧，仅使用该帧加入之前已存在的高斯进行渲染（不泄露未来信息），提取 backbone 特征并完成逐点上下文构造。

采集产物以高斯为核心单位组织：每个训练样本对应一个高斯 seed，包含其上下文特征向量（当前帧与历史帧的原始及渲染特征、观察方向、可见性掩码、逆深度、法线等共 15 类张量）作为输入，以及收敛后的高斯属性（尺度、旋转、不透明度、SH 颜色）作为标签。

\subsubsection{尺度-旋转歧义消解}

2DGS surfel 存在固有的尺度-旋转歧义：交换两个切线方向的尺度并相应旋转 $90^\circ$ 可得到等价表示。若不消解此歧义，训练标签中同一几何形态会对应两种不同的参数组合，导致网络学习不稳定。

数据准备阶段显式消解此歧义：对每个 GT 高斯，强制 $s_1 \geq s_2$；当需要交换时，同步对四元数施加 $90^\circ$ 旋转：$\mathbf{q}' = \text{normalize}((w-z, x+y, y-x, w+z) / \sqrt{2})$。

\subsubsection{训练目标}

网络以多任务损失联合训练，涵盖尺度、旋转、颜色和不透明度全部预测属性。尺度损失在 log 世界尺度空间以 L1 距离度量；旋转损失采用四元数余弦距离 $\mathcal{L}_{\text{rot}} = 1 - |\langle \mathbf{q}_{\text{pred}}, \mathbf{q}_{\text{gt}} \rangle|$；颜色损失对 DC 和高阶 SH 系数分别以 MSE 监督；不透明度损失为 L1 距离。训练采用 Adam 优化器，batch size 为 8192（逐点级），配合梯度裁剪。

\subsection{在线优化}

前馈网络预测的新高斯属性作为初始值插入地图后，仍需经过在线优化进一步收敛。优化以多视角光度损失（L1 与 SSIM 的加权组合）和 2DGS 几何约束损失（深度畸变 + 法线一致性）联合驱动。系统采用多优先级的训练视图选择策略：最新滑窗关键帧具有最高优先级，其次为高损失帧，其余帧通过随机采样补充。优化器采用 \texttt{SparseGaussianAdam}，每次仅更新当前渲染视图下可见的高斯子集，以提升大规模地图的优化效率。
